\section{Discussion and Conclusion}
\label{sec:discussion}

\subsection{Mechanism of the learned head's high-SNR contribution}
\label{ssec:disc_head}

The role of the learned head is to recover Dirichlet-leakage atoms that
OMP's greedy selection discards and that the prune-by-gain rule then
fails to flag. Under~\eqref{eq:dirichlet}, each path produces a
main-lobe atom at the nearest integer Doppler bin and two side-lobes
at $\pm 1$ Doppler offsets carrying $5$--$20\%$ of main-lobe energy
depending on the fractional residual. OMP always picks the main lobe,
but picks a side-lobe only when the side-lobe matched-filter response
beats every other unselected atom --- a function of both main-lobe
amplitude and noise realisation. At low SNR the noise dominates the
selection and the side-lobe is drowned; at high SNR the noise
background falls below the side-lobe amplitudes and the side-lobes
appear at \emph{predictable} integer offsets from active atoms, which
is exactly the pair-feature $\mathbf{v}_{i,s}$ of \eqref{eq:pairfeature}.
The $-37.7\%/-40.2\%/-37.6\%$ pooled paired-BER reductions at
$\snr\in\{20,25,30\}$~dB are therefore not a generic noise-averaging
effect but the signature of a Dirichlet side-lobe support-repair
mechanism that becomes exploitable once the structural predictability
of the side-lobes rises above the residual noise floor. The
attribution of Table~\ref{tab:attribution} (\S\ref{ssec:exp_attribution})
makes this measurable: $>99\%$ of head-added atoms land on
Dirichlet side-lobes / virtual taps.

\subsection{Loss-form and head-capacity ablations}
\label{ssec:disc_lossform}
Two ablations rule out the train-deploy objective gap and a
head-capacity bottleneck. (i) Five support-refinement-head retrainings under a
ListNet-style listwise cross-entropy produce pixel-identical bit
decisions on a shared evaluation cache at $\snr=20$~dB (64 bit errors /
204\,800 bits per seed); the fixed-budget selector reads only the rank order,
so training-time stochasticity does not propagate to deployment.
(ii) Widening pair-MLP and scorer from 16 to 64 takes the head from
1\,041 to 13\,377 parameters ($12.8\times$) and fits BCE
$\approx 1.9\times$ more tightly but yields a non-significant $-5.1\%$
paired reduction at $\snr=20$~dB (McNemar $p=0.34$). Neither
loss-form nor capacity changes the residual BER above the
physical-support oracle floor materially; further improvement requires
an architectural change (non-rank deployment selector), not a training
change.

\subsection{Justification of the 4\,B operating point}
\label{ssec:disc_oversampling}
The $\fs=4B$ operating point is non-standard for AFDM literature
(Bemani 2023~\cite{bemani2023afdm} uses $\fs=B$; most OTFS/AFDM-UWA
papers use $\fs=2B$ Nyquist). The sensitivity test
(\S\ref{ssec:exp_fs}) confirms that $2B$ is $2.4\times$ \emph{worse}
than $4B$ for end-to-end neural receivers under matched hyperparameters,
because $4B$ reduces $k_{\max}$ from 4 to 2 (finer integer Doppler grid,
smaller residual fractional offset, tighter Dirichlet side-lobe energy),
even though $l_{\max}$ doubles and the total dictionary $Q$ grows from
720 to 800. The $4B$ oversampling is therefore not a post hoc
artefact but a defensible design choice for suppressing Doppler ambiguity,
consistent with the oversampled AFDM channel estimator of
Ma~\emph{et al.}~\cite{ma2024oversampledafdm}.

\subsection{Position vs.\ prior AFDM and UWA work}
\label{ssec:disc_priorart}

Table~\ref{tab:priorart} consolidates the design provenance. The
present receiver inherits the sparse-recovery channel-estimation
skeleton from~\cite{benzine2023sparseafdm}, the iterative joint
estimation--detection structure from~\cite{qing2024spafdm,zhao2024gifreeafdm},
the fractional-Doppler Dirichlet expansion from~\cite{yin2024drceafdm},
the $4B$ operating point from~\cite{ma2024oversampledafdm}, the MSML
treatment (first-order) from~\cite{zhou2024msmlafdm}, and the
deep-unfolding template from~\cite{monga2021unrolling,gregor2010lista,borgerding2017amp}.
The 16-pilot reproduction of the SP-AFDM receiver
of~\cite{qing2024spafdm} at its native operating point
matches the published Fig.~3 curve to $\pm 0.2$ log-decade at $4/6$ SNR
cells on a 1\,000-frame budget (Appendix~\ref{app:cao_baseline}). A
direct head-to-head BER comparison is precluded by Zheng's pilot-feasibility
constraint $M(Q{+}1)<\Nsub{-}Q$, structurally incompatible with our
$Q\approx\Nsub/3$ at $\Nsub\geq 1024$.

\begin{table}[t]
\centering
\caption{Structural comparison. SP=superimposed pilots, EP=embedded,
MSML=multi-scale multi-lag, ``stat.\ disc.''=paired-seed significance
testing. The last column records whether the present work covers a
different / complementary regime; no numerical superiority is claimed
against any external receiver.}
\label{tab:priorart}
\footnotesize
\resizebox{\linewidth}{!}{%
\begin{tabular}{lllllll}
\toprule
Work & Waveform & Pilot & Channel & Iter./learned & Stat. & vs.\ this work \\
\midrule
\cite{bemani2023afdm} & AFDM & EP & 2-path, integer Doppler & No/no & no & diff.\ axis \\
\cite{benzine2023sparseafdm} & AFDM & EP & Sparse DD & OMP/no & no & diff.\ axis \\
\cite{qing2024spafdm} & AFDM & SP & Sparse DD & Iter./no & no & complementary; Fig 3 reproduced \\
\cite{zhao2024gifreeafdm} & AFDM & GI-free & Sparse DD & Iter.\ IC/no & no & complementary \\
\cite{yin2024drceafdm} & MIMO-AFDM & Pilot & Frac.\ Doppler & Diag.\ recon./no & no & complementary \\
\cite{ma2024oversampledafdm} & Oversampled AFDM & Pilot & Sparse DD & Iter./no & no & complementary; phys.\ motivation \\
\cite{zhou2024msmlafdm} & AFDM & Pilot & MSML & AFDM-OMP/no & no & stricter physics \\
\cite{afdm2025jointcesd} & AFDM & GI-free & Sparse DD & GI-iter.\ or DNN & no & diff.\ axis (no SP, no UWA) \\
\textbf{This} & AFDM & \textbf{SP} & \textbf{Sparse DD + 1st-MSML} & \textbf{OMP-once + LS/MP-BP + learned} & \textbf{yes} & --- \\
\bottomrule
\end{tabular}
}
\end{table}

\subsection{Limitations}
\label{ssec:limitations}

Five limitations delimit the design statement. (i) The learned head is
deployed as a single trained model; the 3-train-seed envelope at
the $\fs=4B$ hard-decision deployment point on a $3{\times}3$ grid gives
$-5.17\%\pm 0.57$~pp at $\snr=20$~dB, transferred to the headline
recipe by the rank-invariance bridge (\S\ref{ssec:multiseed}). A
strictly empirical $2B+\text{soft}$ retraining envelope is the obvious
follow-up. (ii) Complexity accounting is in per-frame MAC; embedded
latency, memory-hierarchy, and energy on DSP/FPGA are out of scope.
(iii) The simulator is block-fading with first-order MSML; full
broadband per-path waveform resampling and per-frame Doppler jitter
(\cite{zhou2024msmlafdm}) are not modelled. (iv) Head-to-head reproduction
of~\cite{qing2024spafdm} under matched channel statistics is precluded by
the pilot-feasibility constraint $M(Q{+}1)<\Nsub{-}Q$.
(v) Scope is fixed at single-antenna, QPSK, $\rhop=0.5$,
$v_{\max}=3$~m/s; MIMO, higher-order QAM, $\rhop$ sweep, and
high-velocity (>5~m/s) extensions are natural follow-ups.

\subsection{Conclusion}
\label{sec:conclusion}

This paper develops an SP-AFDM joint receiver as a controlled instance
in which the three central methodological gaps in existing
learned receivers can be closed simultaneously: initialization-faithful
unrolling (Claim~\ref{thm:initfaithful}), DA-LS-inclusive
complexity accounting (\S\ref{ssec:complexity_results}), and same-budget
classical heuristic comparison (\S\ref{ssec:exp_omppk}). On the
methodology side, a 1\,041-parameter learned support-refinement head
inserted between the single OMP pass and the iterative LS/MP-BP loop
delivers a consistent $-37.7\%/-40.2\%/-37.6\%$ paired BER reduction
at $\snr\in\{20,25,30\}$~dB under continuity-corrected paired McNemar
($p<10^{-6}$), at $\approx 8\%$ analytical MAC overhead and measured GPU
wall-clock parity. On the mechanism side, the population-level
per-candidate attribution shows $>99\%$ of head additions land on
Dirichlet side-lobes / virtual taps of active physical paths, and the
top-$\Ppath$ physical-support oracle confirms that the learned head
closes $61\%/81\%/98\%$ of the available BER reduction ceiling.

Four directions sharpen the contribution: (a) a strict
$2B+\text{soft}$ retraining-seed envelope to confirm the rank-invariance
bridge empirically; (b) a head-to-head reproduction of~\cite{qing2024spafdm}
at a downscaled operating point; (c) a full broadband-MSML simulator
extension following~\cite{zhou2024msmlafdm}; and (d) an
engineering-deployability study measuring runtime, memory, and energy
on a reference embedded target. None of the four invalidates the
methodology statement of the present paper; they extend its scope.
